\begin{itemize}
    \item \texttt{flip}은 수열의 모든 원소를 뒤집는 연산이다. 즉, 각 $i$에 대해 현재 값 $A_i$를 $1-A_i$로 바꾼다.
    \item \texttt{reverse}는 수열의 순서를 뒤집는 연산이다. 즉, 연산 후의 수열 $B$는 각 $i$에 대해 $B_i = A_{N+1-i}$를 만족한다.
    \item \texttt{flip 후 reverse}는 먼저 모든 원소를 뒤집은 뒤, 그 결과 수열의 순서를 뒤집는 연산이다. 즉, 연산 전 수열이 $A$이고 연산 후 수열이 $B$라면 각 $i$에 대해 $B_i = 1 - A_{N+1-i}$이다.
    \item 채점기의 변형은 오직 $10$번째, $20$번째, $30$번째, $\dots$ 질문에 답한 직후에만 일어날 수 있다. 최종 답안 출력 \texttt{! }$A_1$ $A_2$ $\dots$ $A_N$ 은 질문으로 세지지 않는다. 따라서 예를 들어 $9$번째 질문까지 한 뒤 곧바로 정답을 출력한다면, 그 사이에 수열이 추가로 변하는 일은 없다.
    \item 채점기가 어떤 연산을 적용했는지는 참가자에게 알려지지 않는다.
    \item 최종적으로 출력해야 하는 것은 초기 수열이 아니라, 정답을 출력하는 시점의 현재 수열이다.
\end{itemize}